\documentclass[openany]{ctexbook}
\usepackage{amsmath,amssymb}
\usepackage[hidelinks]{hyperref}
\newcommand{\note}[1]{\footnote{本批译注：#1}}
\title{代数几何基础第五卷预备笔记\\8.14至8.15：校订续译}
\author{据用户上传英译编辑稿译出}
\date{2026年11月15日}
\begin{document}
\maketitle
\chapter*{范围及校订状态}
本批翻译底本第32页全部内容，从“结合8.13、8.5和8.6”起，包含推论8.14、备注8.15及末尾题为“Remorse”的补记。下一处是第33页命题8.16标题。本批为独立增补，未并入原对话合订本。

8.8缺失条件、8.9证明缺口及8.13记号和条件的疑点仍保留；本文不声称已补齐或独立证明这些结果。可确定的排印问题在正文订正并注明原样，不确定内容仅标待核。
\chapter*{8.14\quad 推论（原稿陈述）}
结合8.13、8.5和8.6，按通常方式可得如下推论。\note{此为原作者的引导句；底本未另列完整推导，本译不把它当作已经进行独立验证的声明。}

假设 $f\colon X\to P$ 非分歧，且 $X$ 没有孤立点。〔字迹不清〕$n$。\footnote{原脚注49：“询问 A.G.”。}\note{假设末尾的残字 $n$ 及缺文照录，无法确定其数学含义，不据此补造附加假设。}\note{“直纹的”是原稿的暂用术语，不能径同于通常直纹簇概念；底本8.17专门指出二者不一致。8.13及其依赖的缺文尚未恢复，因此本批完整保存原陈述的充要方向，但不为其提供未经核实的补证。}
\begin{enumerate}
\item[(a)] 假设 $X$ 在 $k$ 上可分。\note{底本用词为 \emph{separable}。此处按原用语译为“可分”，不擅自改成“光滑”或“几何正规”；该旧用语在全书中的精确定义仍须结合前文核定。}
为使那些满足 $\phi_\xi$ 为 $X_{k(\xi)}$-正则、且 $Y_\xi$ 可分的点 $\xi\in P^\vee$ 所成集合，其补集的余维至少为 $2$，充要条件是：$X$ 的每个一维不可约分支 $X_i$ 都由 $X$ 的光滑点组成，并且相对于 $f$ 是“直纹的”；此外，对于 $X$ 的每个满足 $\dim_x X\ge2$ 的闭点 $x$，都有
\note{底本将 $\phi_\xi$ 误排成 $\phi_x i$，按本节一贯记号订正。深度条件印作 $\operatorname{prof}_x X\ge2\operatorname{prof}_x X$，重复了式末的 $\operatorname{prof}_x X$；正文删除该重复。依据是本句随后关于几何正规概形的说明，以及(b)的平行深度条件。按字面把它解释成乘积会使正深度点不满足条件，与原文紧随的说明不相容。此为排印校订，不是宣称已验证整个充要条件。$\operatorname{prof}$ 表示深度。}
\[
\operatorname{prof}_x X\ge2.
\]

（若 $X$ 几何正规，且所有不可约分支的维数均至少为 $2$，这些条件便自动满足。）
\item[(b)] 假设 $X$ 几何正规。为使那些满足 $\phi_\xi$ 为 $X_{k(\xi)}$-正则、且 $Y_\xi$ 几何正规的点 $\xi\in P^\vee$ 所成集合，其补集的余维至少为 $2$，充要条件是：$X$ 的每个维数不超过 $2$ 的不可约分支 $X_i$ 都由 $X$ 的光滑点组成，并且相对于 $f$ 是“直纹的”；此外，对于 $X$ 的每个满足 $\dim_x X\ge3$ 的闭点 $x$，都有
\[
\operatorname{prof}_x X\ge3.
\]
\end{enumerate}

\section*{8.15\quad 备注}
在8.6、8.13和8.14中，我们对 $X$ 作出 $(S_k)$（相应地，$(R_k)$；又相应地，可分、几何正规）这一假设，而希望将同一性质作为超平面截面的结论；唯一可能的例外，是 $\xi$ 属于某个余维至少为 $2$ 的例外集合。

这并不限制一般性；说实话，最好去掉这一预先假设。因为借助第3.4段和5.12的结果，可以立即看出：若 $X$ 不满足所讨论的假设，那么（由第5段），如果存在一个使该假设失效的闭点 $x$，则对于每个满足以下条件的 $\xi$：
\begin{enumerate}
\item $\phi_\xi$ 为 $X_{k(\xi)}$-正则（这个条件只排除一个余维为〔字迹不清〕的集合）；
\item $x\in Y_\xi$（这个条件描述一个余维恰为 $1$ 的集合），
\end{enumerate}
$Y_\xi$ 在 $x$ 处都不满足上述假设。于是，例外集合 $Z\subset P^\vee$ 的余维是 $1$，而不是 $2$。

（对于 $E_k$ 的情形，我也许说得有些过头了；在那里仍需要某种条件：$(S_1)$，也许还需要等维性，也许……）\note{这一保留意见为原文的一部分，必须与前面“不限制一般性”的断言一并阅读。底本的 $E_k$ 记号及缺失的余维数字均未恢复；不能把本段改写成无条件成立的去假设定理。}
\section*{补记：重新考虑（Remorse）}
在8.13和8.14中，只须假设 $f\colon X\to P$ 在 $X$ 的光滑点处非分歧；对于充分性部分，更只须它在 $X$ 的某个开子集 $X'$ 上非分歧，且其补集的余维至少为 $k+1$。
\note{底本将这段另题为“Remorse”，这里译作作者重新考虑前述假设的补记。它没有展开验证，且 $k$ 在8.14不同性质中的取值未在此重述；本译保留原范围和充分性限定，不据此全面替换前面条目的假设。}
\end{document}
